% Generated by scripts/blueprint_gen.py; do not edit.
\chapter{Sheaves on a compact space}\label{chap:S}
\textbf{Setting for all of §3.} X a compact Hausdorff space, \texttt{X : Type} (conventions §7). Sheaves on \texttt{Opens X} with the standard topology. R a sheaf of commutative rings, modules are Mathlib \texttt{SheafOfModules} ✓ (abelian ✓ \texttt{SheafOfModules.instAbelian}, Grothendieck abelian ✓, hence enough injectives ✓ \texttt{IsGrothendieckAbelian.enoughInjectives}).

\begin{theorem}[S1 — Fine sheaves of rings]
\label{S1}
\lean{Cohesive.SupportedIn, Cohesive.IsFine, Cohesive.existsUnique_extend, Cohesive.exists_hom_extend}
\leanok
Let X be a topological space and R a sheaf of rings on X. A global section ρ ∈ R(X) is \emph{supported in} an open U if there is an open V with U ∪ V = X and ρ|\textsubscript{V} = 0 (equivalently supp ρ ⊆ U). R is \emph{fine} if for every finite open cover (U\textsubscript{i}) of X there are ρ\textsubscript{i} ∈ R(X) with Σ\textsubscript{i} ρ\textsubscript{i} = 1 and each ρ\textsubscript{i} supported in U\textsubscript{i}.

\begin{itemize}
\item \textbf{Extension by zero.} If ρ is supported in U, M is a sheaf of R-modules and t ∈ M(U), there is a unique s ∈ M(X) with s|\textsubscript{U} = ρ|\textsubscript{U} · t and s|\textsubscript{V} = 0 on every open V where ρ|\textsubscript{V} = 0.
\item \textbf{For morphisms.} If R is commutative and ρ is supported in U, every morphism φ : M|\textsubscript{U} → N|\textsubscript{U} gives a morphism ψ : M → N with ψ = ρ·φ over opens V ⊆ U and ψ = 0 over opens where ρ vanishes.
\end{itemize}
Deviation: only finite covers are required (the usual notion on compact spaces).
\par\smallskip{\footnotesize Source: blueprint.md §3.}
\end{theorem}
\begin{proof}
\leanok
See the Lean proof and the cited items.
\paragraph{Blueprint notes.}
S is \emph{fine} if for every finite open cover (U\textsubscript{i}) of X there are ρ\textsubscript{i} ∈ S(X) with Σρ\textsubscript{i} = 1 and supp ρ\textsubscript{i} ⊆ U\textsubscript{i} (support = closed set of points where the germ is nonzero). ✗ in Mathlib.

\begin{itemize}
\item \textbf{S1.a Extension by zero of sections.} If supp ρ ⊆ U and t ∈ M(U) for an S-module M, then ρt extends uniquely to a global section vanishing off supp ρ.
\item Finite covers suffice because X is compact; \warn{} [A] uses softness and arbitrary covers.
\end{itemize}
\end{proof}

\begin{theorem}[S2 — Γ is exact on modules over a fine sheaf]
\label{S2}
\lean{Cohesive.exists_local_preimage_of_epi, Cohesive.surjective_globalSections_of_epi}
\leanok
\uses{S1}
Let R be a sheaf of rings on a topological space X and f : M → N an epimorphism of sheaves of R-modules.

\begin{itemize}
\item f is locally surjective: for every s ∈ N(U) and x ∈ U there are an open x ∈ V ⊆ U and t ∈ M(V) with f(t) = s|\textsubscript{V}.
\item If X is compact and R is fine, then Γ(X, f) : M(X) → N(X) is surjective.
\end{itemize}
\par\smallskip{\footnotesize Source: blueprint.md §3.}
\end{theorem}
\begin{proof}
\leanok
See the Lean proof and the cited items.
\paragraph{Blueprint notes.}
For S fine, M ↠ N an epimorphism of S-modules, Γ(X, M) → Γ(X, N) is surjective. Proof: lift locally on a finite cover, glue with a partition of unity via S1.a. Lean: epi ⇔ locally surjective (?) (\texttt{Sheaf.isLocallySurjective\_iff\_epi'} ✓ for sheaves of types/abelian groups; check the module version).
\end{proof}

\begin{theorem}[S3 — Acyclic complexes of S-modules]
\label{S3}
\lean{Cohesive.acyclic_globalSections, Cohesive.exactAt_globalSections, Cohesive.quasiIso_globalSections}
\leanok
\uses{S2}
Let X be compact and R a fine sheaf of rings on X. For ℤ-indexed cochain complexes of sheaves of R-modules (with R-linear differentials):

\begin{itemize}
\item if K is acyclic, the complex of global sections Γ(X, K) is acyclic;
\item degreewise: if K is exact at n, then Γ(X, K) is exact at n;
\item if f : K → L is a quasi-isomorphism, so is Γ(X, f).
\end{itemize}
No boundedness is assumed.
\par\smallskip{\footnotesize Source: blueprint.md §3.}
\end{theorem}
\begin{proof}
\leanok
See the Lean proof and the cited items.
\paragraph{Blueprint notes.}
For S fine, an acyclic complex of S-modules has acyclic complex of global sections. Proof: split into short exact sequences of cycles; apply S2. No sheaf cohomology needed. Degreewise form: exactness of K at n gives exactness of Γ(K) at n (used in S10). Replaces [A, Lemma 7.2.8] for complexes \emph{of} S-modules (S-linear differentials); for complexes whose terms are S-modules but whose differentials are not S-linear (C4), use S3.b.
\end{proof}

\begin{theorem}[S3.b — Soft sheaves]
\label{S3.b}
\lean{Cohesive.IsSoft, Cohesive.IsSoft.of_action, Cohesive.isSoft_of_isFine, Cohesive.surjective_globalSections_of_isSoft, Cohesive.IsSoft.of_shortExact, Cohesive.acyclic_globalSections_of_isSoft}
\leanok
Let R be a sheaf of rings on a topological space X. A sheaf M of R-modules is \emph{soft} if for every closed K ⊆ X, open U ⊇ K and s ∈ M(U) there are t ∈ M(X) and an open V with K ⊆ V ⊆ U and t|\textsubscript{V} = s|\textsubscript{V}.

\begin{itemize}
\item (i) Let S be a fine sheaf of rings on X acting on M through an additive map S(X) → End(M) with 1 acting as the identity and ρ acting as 0 over every open where ρ vanishes. Then M is soft. In particular every module over a fine commutative sheaf of rings is soft.
\item Let X be compact Hausdorff and 0 → A → B → C → 0 a short exact sequence of R-modules with A soft. (ii) Γ(X, B) → Γ(X, C) is surjective. (iii) If B is also soft, so is C.
\item (iv) On compact Hausdorff X, a bounded-below acyclic complex of R-modules with soft terms has acyclic complex of global sections.
\end{itemize}
Deviation: softness is defined through open neighbourhoods of closed sets rather than stalks or sections over closed sets; on compact Hausdorff spaces this is the usual notion.
\par\smallskip{\footnotesize Source: blueprint.md §3.}
\end{theorem}
\begin{proof}
\leanok
\uses{S2}
See the Lean proof and the cited items.
\paragraph{Blueprint notes.}
X compact Hausdorff. Stalk-free softness: a sheaf F of abelian groups is \emph{soft} if every section over an open neighbourhood of a closed K agrees near K with a global section.

\begin{itemize}
\item \textbf{S3.b(i)} A sheaf of abelian groups on which the global sections of a fine sheaf of rings S act (by sheaf endomorphisms, additively, unitally, with ρ· = 0 on opens where ρ vanishes) is soft. In particular every S-module is soft.
\item \textbf{S3.b(ii)} If 0 → A → B → C → 0 is exact and A is soft, Γ(B) → Γ(C) is surjective. Proof: finite cover with local lifts, shrink to closed sets (normality), glue inductively, correcting by global sections of A (softness).
\item \textbf{S3.b(iii)} If A and B are soft, so is C.
\item \textbf{S3.b(iv)} A bounded-below acyclic complex of sheaves (of R-modules) with soft terms has acyclic complex of global sections: the cycle sheaves are soft by induction from the bottom ((iii)), and (ii) lifts global cycles.
\end{itemize}
Replaces [A, Lemma 7.2.8] (Γ-acyclic terms, Stacks 015E) in C4 without sheaf cohomology.
\end{proof}

\begin{theorem}[S4 — Projective modules over a fine sheaf]
\label{S4}
\lean{Cohesive.projective_free, Cohesive.projective_of_isFinFreeSummand, Cohesive.isFinFreeSummand_of_locally, Cohesive.finite_globalSections, Cohesive.projective_globalSections, Cohesive.bijective_globalSections_map}
\leanok
\uses{S2}
Let R be a sheaf of rings on a topological space X. A sheaf of R-modules is a \emph{finite free summand} if it is a retract of R\textsuperscript{m} for some m ∈ ℕ.

\begin{itemize}
\item If X is compact and R is fine, every free R-module R\textsuperscript{(I)} (any index type I) is projective; hence every finite free summand is projective.
\item If X is compact and R is fine and commutative, a module that is locally a finite free summand (each point has an open U with M|\textsubscript{U} a finite free summand) is globally one.
\item For any R and any finite free summand M: Γ(X, M) is a finitely generated projective R(X)-module, and Γ : Hom\textsubscript{R}(M, N) → Hom\textsubscript{R(X)}(Γ(X, M), Γ(X, N)) is bijective for every R-module N.
\end{itemize}
\par\smallskip{\footnotesize Source: blueprint.md §3.}
\end{theorem}
\begin{proof}
\leanok
\uses{S8}
See the Lean proof and the cited items.
\paragraph{Blueprint notes.}
For S fine:

\begin{itemize}
\item \textbf{S4.a} Finite free S-modules are projective (Hom(S, M) = M(X) and S2).
\item \textbf{S4.b} A locally finitely generated, locally projective S-module is a direct summand of a finite free module (finite cover, partition of unity). [A, Lemma 7.2.15, Cor 7.2.13]
\item \textbf{S4.c} Γ is an equivalence between direct summands of finite free S-modules and finitely generated projective S(X)-modules; inverse P ↦ P̃, P̃(U) = S(U) ⊗\textsubscript{S(X)} P. Moreover Hom\textsubscript{S}(P̃, N) = Hom\textsubscript{S(X)}(P, N(X)) for every S-module N. [A, remark after Def 7.2.5, and proof of Lemma 7.2.7]
\item \textbf{S4.d} For such P̃, Γ(P̃ ⊗\textsubscript{S} N) = P ⊗\textsubscript{S(X)} Γ(N). (v1.9: proved in the stronger form "the sectionwise tensor product P̃ ⊗\textsuperscript{pre} N is already a sheaf", on any site and any object, for S a sheaf of commutative rings; needs no fineness.)
\item Lean: \texttt{SheafOfModules.free} ✓, \texttt{IsLocallyFree} ✓.
\end{itemize}
\end{proof}

\begin{theorem}[S5 — Restriction and extension by zero]
\label{S5}
\lean{Cohesive.restrFreeIso, Cohesive.IsFinFreeSummand.restr, Cohesive.extZero, Cohesive.extZero_preservesFiniteLimits, Cohesive.isIso_extZeroAdj_unit, Cohesive.extZero_locally_zero}
\leanok
Let R be a sheaf of rings on a topological space X and U ⊆ X open. Restriction M ↦ M|\textsubscript{U} sends sheaves of R-modules to sheaves of R|\textsubscript{U}-modules on U; it is exact.

\begin{itemize}
\item R\textsuperscript{(I)}|\textsubscript{U} ≅ (R|\textsubscript{U})\textsuperscript{(I)}, and the restriction of a finite free summand is a finite free summand.
\item Extension by zero j\_! is the left adjoint of restriction. The unit M → (j\_! M)|\textsubscript{U} is an isomorphism; every section of j\_! M vanishes near each point outside U; and j\_! is exact.
\item Consequently restriction preserves injective objects.
\end{itemize}
\par\smallskip{\footnotesize Source: blueprint.md §3.}
\end{theorem}
\begin{proof}
\leanok
\uses{S1}
See the Lean proof and the cited items.
\paragraph{Blueprint notes.}
\begin{itemize}
\item \textbf{S5.a} Restriction M ↦ M|U to an open U, exact. Lean: \texttt{overFunctor.\{0\}} ✓ through the project abbreviation \texttt{restr} (conventions §7, decided in P3); exact because it has adjoints on both sides.
\item \textbf{S5.b} Extension by zero j\_! for an open j : U → X: left adjoint to restriction, exact, stalks (j\_!M)\textsubscript{x} = M\textsubscript{x} for x ∈ U and 0 otherwise. ✗ as such in Mathlib, but \texttt{SheafOfModules.pullback} ✓ (left adjoint of \texttt{pushforward} along a continuous functor) applied to the inclusion of opens may already be j\_! (?); exactness still needs a stalk argument.
\item \textbf{S5.c} Restriction preserves injective objects (it has the exact left adjoint j\_!).
\end{itemize}
\end{proof}

\begin{theorem}[S6 — Locality]
\label{S6}
\lean{Cohesive.shortComplex_exact_iff, Cohesive.acyclic_of_locally, Cohesive.quasiIso_of_locally, Cohesive.quasiIso_of_sectionwise}
\leanok
Let R be a sheaf of rings on a topological space X.

\begin{itemize}
\item A short complex M₁ → M₂ → M₃ of R-modules is exact iff every section of M₂ over U killed by the second map comes, near every point of U, from a section of M₁.
\item A complex K of R-modules is acyclic if every point has an open neighbourhood U with K|\textsubscript{U} acyclic; a map f of complexes is a quasi-isomorphism if every point has an open neighbourhood U with f|\textsubscript{U} a quasi-isomorphism.
\item On any site, a map of complexes of sheaves of modules that is a quasi-isomorphism on sections over every object is a quasi-isomorphism.
\end{itemize}
\par\smallskip{\footnotesize Source: blueprint.md §3.}
\end{theorem}
\begin{proof}
\leanok
\uses{S2, S5}
See the Lean proof and the cited items.
\paragraph{Blueprint notes.}
A map of complexes of R-modules is a quasi-isomorphism iff its restrictions to the members of some open cover are; a complex is acyclic iff locally acyclic. Also: a sectionwise homotopy equivalence of presheaves of complexes induces a quasi-isomorphism of the associated sheaf complexes.
\end{proof}

\begin{definition}[S7 — Hom sheaves]
\label{S7}
\lean{Cohesive.homSheaf, Cohesive.homPre, Cohesive.homPost, Cohesive.homSheafComplex, Cohesive.homComplexSectionsEquiv, Cohesive.homSheafOverIso}
\leanok
\uses{S3.b, S5, S6, S9}
Let R be a sheaf of commutative rings on a site C (in particular on X or on the over-site of an open U). For R-modules M, N the \emph{Hom sheaf} 𝓗om(M, N) is the sheaf of R-modules with 𝓗om(M, N)(V) = Hom(M|\textsubscript{V}, N|\textsubscript{V}), R acting through the target, functorial in both variables (pre- and post-composition). For ℤ-indexed complexes K, L, the \emph{Hom complex} 𝓗om(K, L) is the complex of sheaves of R-modules with 𝓗om(K, L)\textsuperscript{n} = ∏\textsubscript{p+n=q} 𝓗om(K\textsuperscript{p}, L\textsuperscript{q}) and differential δz = z ∘ d + (−1)\textsuperscript{n+1} d ∘ z. Its sections over V are the degree-n cochains K|\textsubscript{V} → L|\textsubscript{V}, and 𝓗om(M, N)|\textsubscript{V} ≅ 𝓗om(M|\textsubscript{V}, N|\textsubscript{V}).
\par\smallskip{\footnotesize Source: blueprint.md §3.}
\end{definition}
\paragraph{Blueprint notes.}
For complexes M, N of R-modules, 𝓗om(M, N) is the complex of sheaves of R-modules with 𝓗om(M, N)(V) = Hom\textsuperscript{•}(M|V, N|V). ✗ for \texttt{SheafOfModules} (Mathlib has \texttt{presheafHom} ✓ for type-valued Homs only).

\begin{itemize}
\item \textbf{S7.a′} (v1.9, replaces S7.a) \emph{Termwise source action.} Let a ring T act on each term M\textsuperscript{i} of a complex of R-modules by R-linear sheaf endomorphisms (additively, unitally; not necessarily commuting with the differential). Precomposition gives an action of T on each term 𝓗om(M, N)\textsuperscript{n} by R-linear sheaf endomorphisms, additive and unital, and vanishing on V whenever the action on every M\textsuperscript{i} vanishes on V. It is not a chain map. With T = S(X) for a fine S this is the input of S3.b(i) (\texttt{IsSoft.of\_action}), used in C4 step 2. (The old S7.a, "a complex of S-modules via the source", is false as stated when the differentials of M are not S-linear, which is exactly the case in C3/C4.)
\item \textbf{S7.b} If W is strictly perfect (S9), 𝓗om(W, −) preserves quasi-isomorphisms.
\item \textbf{S7.c} If I is a bounded-below complex of injectives, then 𝓗om(C, I) is acyclic for every acyclic C (sections over V use S5.c and \texttt{CochainComplex.isKInjective\_of\_injective} ✓).
\end{itemize}

\begin{definition}[S8 — Tensor products and K-flat resolutions]
\label{S8}
\lean{Cohesive.tensorBifunctor, Cohesive.tensorC, Cohesive.tensorCMap₂, Cohesive.IsFlat, Cohesive.IsKFlat, Cohesive.restrTensorIso}
\leanok
\uses{S5, S9}
Let R be a sheaf of commutative rings on a site.

\begin{itemize}
\item The \emph{tensor product} M ⊗ N of sheaves of R-modules is the sheafification of the presheaf V ↦ M(V) ⊗\textsubscript{R(V)} N(V).
\item The \emph{tensor product of complexes} K ⊗ L is the total complex of the double complex K\textsuperscript{p} ⊗ L\textsuperscript{q} (Koszul sign on the differential of the second factor); it is functorial in both variables.
\item M is \emph{flat} if M ⊗ − preserves monomorphisms; a complex K is \emph{K-flat} if K ⊗ L is acyclic for every acyclic complex L.
\item On a space X, tensor products commute with restriction to an open U: (M ⊗ N)|\textsubscript{U} ≅ M|\textsubscript{U} ⊗ N|\textsubscript{U}.
\end{itemize}
\par\smallskip{\footnotesize Source: blueprint.md §3.}
\end{definition}
\paragraph{Blueprint notes.}
This package replaces Malgrange's flatness theorem (feasibility report, G6).

\begin{itemize}
\item \textbf{S8.a} Tensor product and base change of sheaves of modules along R → S. Lean: \texttt{PresheafOfModules} monoidal ✓ plus sheafification ✓; no monoidal structure on \texttt{SheafOfModules} ✗.
\item \textbf{S8.b} Flat R-modules; j\_! R|\textsubscript{U} is flat (stalks, S5.b).
\item \textbf{S8.c} Every R-module is a quotient of a direct sum of modules j\_! R|\textsubscript{U}.
\item \textbf{S8.d} Every bounded-above complex has a quasi-isomorphism from a bounded-above complex of flat modules; bounded-above complexes of flats are K-flat. [Stacks, chapter \emph{Cohomology of Sheaves}, section on K-flat complexes; look up the exact tags in phase 2].
\item \textbf{S8.e} If K is K-flat and acyclic then K ⊗ L is acyclic for every L (resolve L by S8.d).
\item \textbf{S8.f} Restriction to opens preserves K-flatness (use j\_!).
\item \textbf{S8.g} If K, K′ are K-flat and connected by a roof of quasi-isomorphisms, so are S ⊗\textsubscript{R} K and S ⊗\textsubscript{R} K′ (resolve the middle term by S8.d). (v1.9) Lean form: for a single quasi-isomorphism K → K′ between K-flat complexes. For a roof K ← L → K′ with K′ bounded (e.g. strictly perfect) and L arbitrary, as in \texttt{IsPerfect}: first replace L by τ\textsubscript{≤b}L (\texttt{CochainComplex.truncLE}; L → K′ a quasi-isomorphism with K′ in degrees ≤ b makes τ\textsubscript{≤b}L → L one), then resolve τ\textsubscript{≤b}L by S8.d and apply S8.g to both legs.
\end{itemize}

\begin{theorem}[S9 — Perfect complexes]
\label{S9}
\lean{Cohesive.IsStrictlyPerfect, Cohesive.IsPerfect, Cohesive.perf_Q_obj_iff, Cohesive.IsPerfectIn, Cohesive.IsPerfect.exists_isPerfectIn, Cohesive.IsPerfectIn.exactAt}
\leanok
Let R be a sheaf of rings on a topological space X; complexes are ℤ-indexed cochain complexes of sheaves of R-modules.

\begin{itemize}
\item K is \emph{strictly perfect} if K\textsuperscript{n} = 0 outside some interval [a, b] and every K\textsuperscript{n} is a finite free summand (on any site).
\item K is \emph{perfect} if every point has an open neighbourhood U, a strictly perfect complex P of R|\textsubscript{U}-modules and quasi-isomorphisms K|\textsubscript{U} ← L → P.
\item D\textsubscript{perf}(X, R) is the full subcategory of the derived category D(R) on the objects isomorphic to Q(K) for some perfect K; Q(K) lies in D\textsubscript{perf}(X, R) iff K is perfect.
\item If X is compact, every perfect K has local models as above all concentrated in one interval [a, b], and then the homology sheaves of K vanish outside [a, b].
\end{itemize}
Deviation: perfectness is a roof of quasi-isomorphisms on each U instead of an isomorphism in D(R|\textsubscript{U}); the two are equivalent.
\par\smallskip{\footnotesize Source: blueprint.md §3.}
\end{theorem}
\begin{proof}
\leanok
\uses{S5}
See the Lean proof and the cited items.
\paragraph{Blueprint notes.}
\begin{itemize}
\item \textbf{S9.a} A complex of R-modules is \emph{strictly perfect} if it is bounded and each term is a direct summand of a finite free module. [A, Def 7.2.11]
\item \textbf{S9.b} A complex K is \emph{perfect} if every point has an open neighbourhood U, a strictly perfect P on U and quasi-isomorphisms K|U ← L → P (\texttt{QuasiIso} ✓). Equivalent to Q(K|U) ≅ Q(P) in D(R|U), but avoids a derived category for every open set. Lean: with \texttt{overFunctor} (chosen in P3), "strictly perfect on U" lives over \texttt{R.over U} on the over-site \texttt{Over U}, not over a ring on \texttt{Opens U}; S9.d and S10 must use the same API.
\item \textbf{S9.c} D\textsubscript{perf}(X, R) is the full subcategory of \texttt{DerivedCategory (SheafOfModules R)} ✓ on objects isomorphic to Q(K) with K perfect (perfectness is invariant under quasi-isomorphism).
\item \textbf{S9.d} On compact X every perfect complex is globally bounded (uniform degree range and number of generators), by a finite cover.
\end{itemize}
\end{proof}

\begin{theorem}[S10 — Globalization over a fine sheaf]
\label{S10}
\lean{Cohesive.exists_strictlyPerfect_quasiIso, Cohesive.nullHomotopic_locally, Cohesive.lift_locally, Cohesive.split_locally, Cohesive.isFinFreeSummand_locally_of_acyclic, Cohesive.homology_finGen_locally, Cohesive.exists_epi_free_of_locally}
\leanok
\uses{S1, S11, S2, S4, S9}
Let X be compact and R a fine sheaf of commutative rings on X. Every perfect complex K of R-modules receives a quasi-isomorphism P → K from a strictly perfect complex P.

Supporting lemmas, for a sheaf of rings R on any site:

\begin{itemize}
\item a chain map from a strictly perfect complex to an acyclic complex is locally null-homotopic;
\item for E strictly perfect and a quasi-isomorphism s : L → K, every chain map E → K lifts locally through s up to homotopy;
\item an epimorphism onto a finite free summand splits locally; in a bounded acyclic complex in degrees [m, n] whose terms other than the (m+1)-st are finite free summands, that term is locally a finite free summand;
\item if P is strictly perfect and exact in all degrees ≥ k, then H\textsuperscript{k−1}(P) is locally a quotient of a finite free module.
\end{itemize}
And on X with R fine: a module that is locally a quotient of a finite free module is a quotient of a finite free module.
\par\smallskip{\footnotesize Source: blueprint.md §3.}
\end{theorem}
\begin{proof}
\leanok
\uses{S7}
See the Lean proof and the cited items.
\paragraph{Blueprint notes.}
For S fine, every globally bounded perfect complex of S-modules is quasi-isomorphic to a strictly perfect complex of S-modules. [A, Lemma 7.2.16, with Lemmas 7.2.9, 7.2.10]

\begin{itemize}
\item Sub-lemmas (v1.8), on any space Y and sheaf of rings R on Y (applied on opens U through Mathlib's \texttt{Opens.sheafOfModulesEquivOver}):
\begin{itemize}
\item \textbf{S10.a Local null-homotopy.} A chain map from a strictly perfect complex to an acyclic complex is locally null-homotopic (finitely many sections, degree by degree, by local surjectivity of epis, S2 local part).
\item \textbf{S10.b Local lifting.} For E strictly perfect and a quasi-isomorphism s : L → K, every chain map E → K lifts locally through s up to homotopy (S10.a on the cone of s).
\item \textbf{S10.c Local splitting.} An epimorphism onto a direct summand of a finite free module splits locally. Hence, for a bounded acyclic complex C in degrees [m, n] whose terms are summands of finite free modules except possibly C\textsuperscript{m+1}: splitting from the top, the cycle sheaves Z\textsuperscript{i} (i ≥ m + 2) are locally such summands, and C\textsuperscript{m+1} ≅ C\textsuperscript{m} ⊕ Z\textsuperscript{m+2} locally, so C\textsuperscript{m+1} is locally a summand of a finite free module too.
\item \textbf{S10.d} The top homology sheaf of a strictly perfect complex is locally finitely generated (S10.c on the part above it).
\item \textbf{S10.e Global generation.} For R fine on compact X, a locally finitely generated R-module is a quotient of a finite free module (finite cover, S1.a).
\end{itemize}
\item Proof outline (v1.8): by S9.d, K has local strictly perfect models in [a, b] and homology in [a, b]. (1) Downward induction from b to a − 1 builds finite free S-modules Fⁿ and F → K, iso on homology above the current degree, as in [TT, Lemma 1.9.5]: the next homology sheaf of the cone is locally the top homology of the cone of a local lift F|U → P\textsubscript{U} (S10.b), a strictly perfect complex, so locally finitely generated (S10.d), hence a quotient of a finite free module (S10.e); lift to global cycles (S4.a). (2) B = coker(F\textsuperscript{a−1} → F\textsuperscript{a}) and P = (B → F\textsuperscript{a+1} → ⋯ → F\textsuperscript{b}), with P → τ\textsubscript{≥a}K ≃ K a quasi-isomorphism. Locally the model P\textsubscript{U} lifts to a quasi-isomorphism P\textsubscript{U} → P|U (S10.b); its cone is bounded and acyclic in degrees [a − 1, b] with terms summands of finite free modules except cone\textsuperscript{a} = P\textsubscript{U}\textsuperscript{a+1} ⊕ B; S10.c makes B locally a summand of a finite free module, and S4.b globally. So P is strictly perfect, hence K-projective (S4.a, S11), and P → τ\textsubscript{≥a}K lifts to P → K (S11).
\item \warn{} The proof of [A, Lemma 7.2.16] applies [A, Lemma 7.2.10] on small opens U, which needs free S|U-modules to be projective (true for soft S on paracompact U, not given by S1 on compact X); replaced by local lifting and local splitting (S10.a–c) and S4.b (v1.8).
\item \warn{} Compactness replaces [A, Lemma 7.2.12] (finite covering dimension).
\end{itemize}
\end{proof}

\begin{theorem}[S11 — Derived-category facts from Mathlib]
\label{S11}
\lean{Cohesive.exists_injective_resolution, Cohesive.exists_lift_of_isKProjective}
\leanok
Let C be an abelian category.

\begin{itemize}
\item If C has enough injectives, every complex K with K\textsuperscript{i} = 0 for i < n has a quasi-isomorphism K → I to a complex I of injective objects with I\textsuperscript{i} = 0 for i < n (so I is K-injective).
\item If P is K-projective and s : L → K is a quasi-isomorphism, every chain map g : P → K lifts through s up to homotopy: there is f : P → L with s ∘ f homotopic to g.
\end{itemize}
Morphisms in D(C) into a K-injective complex, or out of a K-projective one, are homotopy classes of chain maps; bounded complexes of projectives are K-projective.
\par\smallskip{\footnotesize Source: blueprint.md §3.}
\end{theorem}
\begin{proof}
\leanok
See the Lean proof and the cited items.
\paragraph{Blueprint notes.}
\begin{itemize}
\item Bounded-below injective resolutions: \texttt{exists\_quasiIso\_injective} ✓ (\texttt{DerivabilityStructureInjectives}).
\item Hom in D into a K-injective complex equals homotopy classes: \texttt{CochainComplex.IsKInjective} ✓, \texttt{Qh\_map\_bijective\_of\_isKInjective} ✓.
\item Bounded complexes of projectives are K-projective: \texttt{isKProjective\_of\_projective} ✓.
\item Universes: \texttt{HasDerivedCategory.standard} ✓ as a local instance (conventions §7).
\end{itemize}
\end{proof}
